\begin{abstract}
While reproducibility assessment tools operate after experiments run, researchers lack guidance for predicting which datasets will yield reproducible results before writing code. We demonstrate that dataset metadata completeness predicts reproducibility variance in machine learning benchmarks. Analyzing 300 OpenML datasets with 21,312 matched experimental runs, we find that top-quartile metadata completeness predicts a 42.1\% reduction in interquartile range (IQR) of performance outcomes compared to bottom-quartile datasets (95\% CI: 39.1--51.7\%). Mediation analysis identifies preprocessing entropy as the dominant mechanism, mediating 64.7\% of the effect (Sobel $Z = 16.02$, $p < 0.0001$): complete documentation constrains preprocessing choices, reducing implementation heterogeneity and thereby reducing variance. The effect persists in early-run subsamples (90.7\% preservation) and within single algorithm families ($p < 0.001$), ruling out reverse causality and algorithm-mix confounds. Our work reframes reproducibility from a binary property of papers assessed post-hoc to a continuous property of datasets predicted proactively.
\end{abstract}
